\documentclass[10pt,a4paper]{amsart}
\usepackage[margin=19mm]{geometry}
\usepackage{amsmath,amssymb,mathtools}
\usepackage[scr=rsfs,cal=euler]{mathalfa}
\usepackage{xfrac}
\usepackage[unicode,hidelinks,bookmarks=false]{hyperref}
\newcommand{\ie}{i.e.\ }
\newcommand{\cf}{cf.\ }
\newcommand{\resp}{respectively\ }
\newcommand{\Subsection}{\subsection}
\providecommand{\nobreakdash}{\nobreak\hbox{-}}
\setlength{\emergencystretch}{2em}
\multlinegap0pt
\usepackage{bm}
\usepackage{bbm}
\usepackage{dsfont}
\usepackage{xspace}
\usepackage{cancel}
\usepackage{mathtools}
\usepackage{cleveref}
\usepackage{amsthm}
\makeatletter
\@ifundefined{thm}{\newtheorem{thm}{Thm}}{}
\@ifundefined{lemma}{\newtheorem{lemma}[thm]{Lemma}}{}
\@ifundefined{remark}{\newtheorem{remark}[thm]{Remark}}{}
\@ifundefined{theorem}{\newtheorem{theorem}[thm]{Theorem}}{}
\makeatother
\newcommand{\Arr}{\operatorname{Arr}}
\newcommand{\F}{\mathcal{F}}
\newcommand{\NF}{\operatorname{NormFan}}
\newcommand{\NP}{\operatorname{Root}}
\newcommand{\OO}{\mathcal{O}}
\newcommand{\R}{\mathbb{R}}
\newcommand{\vv}{\upsilon}
\title[Root polytopes, flow polytopes, and order polytopes]{Root polytopes, flow polytopes, and order polytopes}
\author{Konstanze Rietsch, Lauren Williams}
\date{Source: arXiv:2406.15803v3 (CC BY 4.0)}
\begin{document}
\maketitle

\noindent\emph{Source and licence.} Root polytopes, flow polytopes, and order polytopes. Version arXiv:2406.15803v3 (CC BY 4.0), distributed under the Creative Commons Attribution 4.0 International licence, as recorded in the arXiv licence metadata for this exact version and shown on the abstract page https://arxiv.org/abs/2406.15803v3. Full rights notice: licenses/arxiv-2406.15803-CC-BY-4.0.txt.\par\smallskip

\noindent\textit{Editorial scope.} Counted unit: the Remark whose source label is \texttt{rem:graded} in the article (source0.tex lines 2410--2418, source label \texttt{rem:graded}), delivered together with the article's own complete original proof of it (source0.tex lines 2421--2473, one \texttt{proof} environment of 53 lines). No mathematics is written, repaired or re-derived: statement and proof are the article's own text line for line. Required context retained uncounted: lem:containstar (lines 1411--1415); lem:rays (lines 2238--2247); thm:refine (lines 2297--2308); lem:orderideal (lines 2356--2364). Every definition, display and label of those lines is retained unchanged. Omitted from this package and not redistributed: 1--1410, 1416--2237, 2248--2296, 2309--2355, 2365--2409, 2419--2420, 2474--3369. Nothing inside a retained excerpt is omitted or altered: every retained body is delivered exactly as the article's lines are written, so no omission receipt of any kind accompanies this package. Citations: the retained excerpts carry no citation command at all, so no bibliography entry of the article is transcribed and no reference is redistributed. Reference adaptation: none was needed and none was made; every reference inside the delivered unit resolves inside it, so no unresolved reference or citation is produced. Printed slips kept literal: no printed word, symbol, hypothesis or number of the counted statement or of its proof was altered. Layout and class: the article's own class is \texttt{amsart} with packages including inputenc, etex, adjustbox, pstricks, amssymb, pict2e, pst-plot, xypic, quiver, tikz, tikz-cd, hyperref, geometry, rotating; the wrapper is the lane's pinned \texttt{amsart} editorial wrapper at 10pt on A4, which restates the theorem-like environments the retained text needs and adds the article's own macro definitions that the retained text needs (\texttt{\textbackslash Arr}, \texttt{\textbackslash F}, \texttt{\textbackslash NF}, \texttt{\textbackslash NP}, \texttt{\textbackslash OO}, \texttt{\textbackslash R}, \texttt{\textbackslash vv}), with extra wrapper packages (bm, bbm, dsfont, xspace, cancel, mathtools, cleveref). The delivered numbering is the unit's own. Assets: no retained span contains a figure environment, a graphics-inclusion command, a PDF-page inclusion, a table or a TikZ picture; no image file is copied into this package, the package declares no asset dependency and no image file is redistributed with it. Licence: version arXiv:2406.15803v3 of this article is distributed under the Creative Commons Attribution 4.0 International licence, as recorded in the arXiv licence metadata for this exact version and shown on the abstract page https://arxiv.org/abs/2406.15803v3. A full rights notice is delivered at licenses/arxiv-2406.15803-CC-BY-4.0.txt. Characters: every retained excerpt is pure ASCII, so the delivered engine, XeLaTeX with its default Latin Modern text font, reports no missing character.
\begin{lemma}\label{lem:containstar}
Let $M:\Arr(Q) \to \R_{\geq -1}$ be a facet arrow-labeling.
Then each facet component of $Q$ (with respect to $M$)
contains some starred vertex.
\end{lemma}

\begin{lemma}\label{lem:rays}
Let $P$ be an arbitrary finite poset.
Then the rays of the 
	face fan $\F_{Q_{\hat{P}}}$ of the root polytope $\NP(Q_{\hat{P}})$
of the starred quiver $Q_{\hat{P}}$
coincide with the rays of the (inner) normal fan $\NF(\OO(P))$ of 
the order polytope $\OO(P)$;  these rays are in bijection 
with arrows of $Q_{\hat{P}}$, or equivalently, with cover relations in 
the Hasse diagram of $\hat{P}$.
\end{lemma}

\begin{theorem}\label{thm:refine}
Let $P=\{\vv_1,\dots,\vv_n\}$ be a finite ranked poset.
	Then the face fan $\F_{Q_{\hat{P}}}$ of the root polytope $\NP(Q_{\hat{P}})$
 of the starred quiver $Q_{\hat{P}}$
refines the (inner) normal fan $\NF(\OO(P))$ of 
the \emph{order polytope} $\mathcal{O}(P)$ of $P$:
the two fans have the same set of rays, and each 
maximal cone of $\NF(\OO(P))$ is a union of maximal cones of 
	$\F_{Q_{\hat{P}}}$.  Moreover, if $P$ is a graded poset, then 
	$\F_{Q_{\hat{P}}}$ coincides with 
 $\NF(\OO(P))$.
\end{theorem}

\begin{lemma} \label{lem:orderideal}
Suppose that $P=\{\vv_1,\dots,\vv_n\}$ is a ranked poset. 
Let $M$ be a facet arrow-labeling of the starred quiver $Q_{\hat{P}}$, and $L$ the corresponding $\bullet$-labeling. Consider the associated facet components  $C_0$ and $C_1$ of $Q_{\hat{P}}$ containing $\star_0$ and $\star_1$, respectively.  

The facet components $C_0$ and $C_1$ are disjoint. The facet component $C_0$ consists of a set of vertices $S\sqcup\{\star_0\}$ along with all arrows between them, where the set $S$ forms a lower order ideal in $P$. For the facet component $C_1$, the set of normal vertices in $C_1$ 
	forms a filter in $P$,
	namely $\{\vv_1,\dots,\vv_n\} \setminus S$. 

\end{lemma}

\begin{remark}\label{rem:graded}
If the poset $P$ in \Cref{lem:orderideal} is not just ranked but also graded,
then the proof in \Cref{lem:orderideal} extends to show that 
the facet component $C_1$ consists of 
the vertices $\{\star_1\} \cup 
\{\vv_1,\dots,\vv_n\} \setminus S $ together with 
all arrows %$\vv_i \to \vv_j$ 
joining two elements in this set.
\end{remark}

\begin{proof}[Proof of \Cref{thm:refine}]
The statement that the rays of the two fans coincide was already proved in \cref{lem:rays}. To show that each maximal cone of $\NF(\OO(P))$ is a union of maximal cones of 
	$\F_{Q_{\hat{P}}}$, it suffices to show that each maximal
	cone of $\F_{Q_{\hat{P}}}$ is contained
in a maximal cone of $\NF(\OO(P))$. Equivalently, we need to show 
	that the rays of each maximal cone of $\F_{Q_{\hat{P}}}$ 
 are a subset of the rays
of some 
 maximal cone of $\NF(\OO(P))$.

	Each maximal cone $c_M$ of $\F_{Q_{\hat{P}}}$ comes from a facet arrow-labeling $M$
of $Q_{\hat{P}}$, and the rays in $c_M$ are indexed
by those 
 arrows $a\in \Arr(Q_{\hat{P}})$ for which 
$M(a) = -1$. These are the arrows appearing in 
the facet components
of $Q_{\hat{P}}$ (with respect to $M$).

By \Cref{lem:containstar}, each facet component of $Q_{\hat{P}}$
contains some starred vertex, so the vertices of $Q_{\hat{P}}$ 
lie in either $C_1$, the facet component of $\star_1$,
or $C_0$, the facet component of $\star_0$.
By \Cref{lem:orderideal}, the vertices of the 
facet component $C_0$ form a lower order ideal of $P$, and the 
vertices of $C_1$ form a filter of $P$.
Therefore each arrow indexing a ray in $c_M$ -- 
which is by definition an arrow appearing in a facet components of $Q_{\hat{P}}$ --
connects either two vertices in the filter $I$, or two vertices
in the complement of the filter $\{\vv_1,\dots,\vv_n\} \setminus I$.
By \Cref{rem:graded},  if $P$ is graded,
then the arrows appearing in $C_i$ (for $i=0$ or $1$)
comprise \emph{all} arrows connecting two vertices in $C_i$.

Meanwhile, each maximal cone $c'_I$ of $\NF(\OO(P))$ 
corresponds to a 
vertex 
 of $\OO(P)$, which is the 
characteristic function $\chi_I$ of a filter $I$ of $P$.
The rays in $c'_I$ correspond to the 
edges $\vv_i \lessdot \vv_j$ of the Hasse diagram of $\hat{P}$,
where $\chi_I(\vv_i) = \chi_I(\vv_j)$.
Therefore each ray of $c'_I$ corresponds to an arrow in $\Arr(Q_{\hat{P}})$
which connects either two vertices in the filter $I$
or two vertices in the complement of the 
filter $\{\vv_1,\dots,\vv_n\} \setminus I$.

This shows that the set of rays in a maximal cone 
$c_M$ 
	of $\F_{Q_{\hat{P}}}$ 
 are a subset of the rays
of a corresponding maximal cone 
 $c'_I$ of $\NF(\OO(P))$, with equality when $P$ is graded.
\end{proof}

\end{document}
